\paragraph{AQ-KO-RANK: a correction to an ancillary dimension claim.}
With Definition 2.5, p.~8, of the same paper, the dimension
sentence in Lemma 4.6, p.~34, requires one additional rank:
for $K=\mathbb Q(\mu_m)$ with $m>2$,
\[
 \operatorname{rank}_{\mathbb Z}\Xi(K)
       =\frac{\varphi(m)}2+1,
 \tag{AQ.KO.1}
\]
where the paper prints $\varphi(m)/2$.
\begin{proof}
Put $r=\varphi(m)/2$ and choose one embedding
$\sigma_j$ from each of the $r$ complex-conjugate pairs.
Since $K$ itself is a CM field, the additional CM-subfield
condition in that definition is automatic.
A function $\phi:\Sigma_K\to\mathbb Z$ belongs to
$\Xi(K)$ precisely when all the sums
$\phi(\sigma_j)+\phi(\overline\sigma_j)$ equal
one common integer $w$.
Consequently
\[
 \Xi(K)\longrightarrow\mathbb Z^{r+1},\qquad
 \phi\longmapsto(w,\phi(\sigma_1),\ldots,\phi(\sigma_r))
\]
is an isomorphism of abelian groups: the inverse sends
$(w,a_1,\ldots,a_r)$ to the function with values
$a_j$ and $w-a_j$ on the $j$th pair.
This proves (AQ.KO.1).  The map
$\phi\mapsto2\phi-w$ used in the cited rank argument
has the full group of constant functions as its kernel,
rather than being injective.  Already $m=3$ gives
$\Xi(\mathbb Q(\mu_3))=\mathbb Z^2$.
\end{proof}
This calculation concerns that rank sentence; the
period-comparison results above are cited with their
own hypotheses, and their complete external proofs
are not certified by this source-scope check.

